\documentclass[11pt,a4paper]{article}
\usepackage[utf8]{inputenc}
\usepackage[T1]{fontenc}
\usepackage[brazil]{babel}
\usepackage[top=2.5cm,bottom=2cm,left=2.5cm,right=2cm]{geometry}
\usepackage{booktabs}
\usepackage{enumitem}
\usepackage{hyperref}

\title{\textbf{Relatório de Acompanhamento Diário -- Experiência Criativa}\\
\large Estudo de Caso: Rede de Farmácias (PUCPR)}
\author{\textbf{Equipe:} Gustavo Kawano, Henrique Papai, Edson Renato}
\date{Setembro de 2026}

\begin{document}

\maketitle

\section{Identificação da Equipe}
\textbf{Participantes:} Gustavo Kawano, Henrique Papai e Edson Renato.

\section{Configuração Obrigatória de Usuários (Menor Privilégio)}
Usuários padrão configurados no ambiente com privilégios mínimos:
\begin{itemize}[noitemsep]
    \item \textbf{Sistema (OS / PostgreSQL / Docker):} \texttt{teste} | \textbf{Senha:} \texttt{T35t3} \\
    \textit{Privilégio:} Apenas consulta (\texttt{SELECT}) a tabelas públicas e view anonimizada. Não-root no container.
    \item \textbf{Aplicações (Keycloak / Web):} \texttt{teste@pucparana.com} | \textbf{Senha:} \texttt{Teste@2026} \\
    \textit{Privilégio:} Papel restrito de \texttt{cliente} no Keycloak, sem acesso administrativo.
\end{itemize}

\section{Backlog}
\begin{itemize}[noitemsep]
    \item Segmentação de rede (DMZ e Interna) com Docker Compose.
    \item WAF ModSecurity com OWASP CRS e regras DLP (bloqueio de CPF e cartão).
    \item IDS/IPS Suricata com 11 regras alinhadas ao MITRE ATT\&CK.
    \item Aplicação e-commerce (Flask), PostgreSQL com view LGPD e Keycloak (MFA).
    \item SIEM Wazuh completo (Manager, Indexer, Dashboard e 2 Agentes).
    \item Script de anonimização LGPD ($k$-anonimato).
    \item Migração para 2 VMs Linux com pfSense e execução dos 4 cenários de teste.
\end{itemize}

\section{Atividades Realizadas}
\begin{itemize}[noitemsep]
    \item \textbf{Infraestrutura:} Criação e subida de \texttt{docker-compose.dmz.yml} e \texttt{docker-compose.internal.yml} (10 containers operacionais).
    \item \textbf{Defesa de Borda:} ModSecurity bloqueando SQL Injection (\texttt{403}) e ferramentas de scan; regras DLP ativas contra vazamento de CPF.
    \item \textbf{Detecção:} Suricata ativo monitorando tráfego lateral e exfiltração.
    \item \textbf{Identidade e Dados:} Keycloak rodando com realm e MFA importados; PostgreSQL com view anonimizada e script \texttt{anonymize.py} validado.
    \item \textbf{SIEM:} Cluster Wazuh Indexer inicializado (\texttt{GREEN}), Dashboard ativo na porta 5601 e agentes DMZ/Interno registrados.
\end{itemize}

\section{Atividades em Desenvolvimento}
\begin{itemize}[noitemsep]
    \item Transferência dos serviços para as 2 VMs Linux definitivas conectadas ao pfSense de laboratório.
    \item Ajuste da interface de captura física do Suricata no Linux.
    \item Automação em scripts (\texttt{.sh}) dos 4 cenários de teste de invasão/defesa.
\end{itemize}

\section{Principais Dificuldades}
\begin{itemize}[noitemsep]
    \item \textbf{Permissões no WAF e Suricata:} Imagens que necessitavam de escrita interna falhavam com montagem \texttt{:ro}. Resolvido injetando as regras no caminho pós-CRS (\texttt{RESPONSE-999-*.conf}).
    \item \textbf{Hardening no Windows/Docker:} Diretiva \texttt{cap\_drop: ALL} impedia a inicialização do PostgreSQL e Wazuh Agent. Resolvido ajustando capacidades mínimas e isolando na rede interna.
    \item \textbf{Inicialização do OpenSearch:} O Wazuh Indexer subiu desconfigurado. Resolvido executando a rotina do \texttt{securityadmin.sh} com certificados TLS internos.
\end{itemize}

\end{document}
